\documentclass[conference]{IEEEtran}
\IEEEoverridecommandlockouts

\usepackage{cite}
\usepackage{amsmath,amssymb,amsfonts}
\usepackage{algorithmic}
\usepackage{graphicx}
\usepackage{textcomp}
\usepackage{xcolor}
\usepackage{hyperref}
\usepackage{booktabs}
\usepackage{multirow}
\usepackage{array}
\usepackage{url}
\usepackage{float}
\usepackage{subcaption}
\usepackage{listings}
\usepackage{tikz}
\usepackage{pgfplots}
\pgfplotsset{compat=1.18}

\hypersetup{
  colorlinks=true,
  linkcolor=blue,
  citecolor=blue,
  urlcolor=blue,
}

\lstset{
  basicstyle=\footnotesize\ttfamily,
  breaklines=true,
  frame=single,
  keywordstyle=\color{blue},
  commentstyle=\color{gray},
  stringstyle=\color{orange},
}

\def\BibTeX{{\rm B\kern-.05em{\sc i\kern-.025em b}\kern-.08em
    T\kern-.1667em\lower.7ex\hbox{E}\kern-.125emX}}

\begin{document}

%=============================================================================
\title{Predicting Global Health Outcomes: A Multi-Source Analytics\\
Pipeline Integrating WHO, World Bank, and Disease APIs}

\author{
  \IEEEauthorblockN{Vecha Jatadhar}
  \IEEEauthorblockA{MSc Data Analytics\\
    National College of Ireland\\
    Dublin, Ireland}
  \and
  \IEEEauthorblockN{Pravalika Revelli}
  \IEEEauthorblockA{X25161261\\
    National College of Ireland\\
    Dublin, Ireland}
  \and
  \IEEEauthorblockN{Marla Srija}
  \IEEEauthorblockA{X24295256\\
    National College of Ireland\\
    Dublin, Ireland}
}

\maketitle

%=============================================================================
\begin{abstract}
Global health disparities remain one of the most pressing challenges of the
21st century, with life expectancy varying by more than 30 years between the
wealthiest and poorest nations. This paper presents an end-to-end analytics
pipeline that integrates heterogeneous data from three open APIs---the World
Health Organisation Global Health Observatory (WHO GHO), the World Bank Open
Data API, and the disease.sh COVID-19 API---to investigate the socioeconomic
and health-system determinants of life expectancy across 185+ countries spanning
2010--2022. Raw semi-structured JSON data is ingested into MongoDB, transformed
via a programmatic ETL pipeline, and stored in a structured PostgreSQL schema.
Three machine learning models are applied: K-Means clustering segments countries
into five distinct health profiles; a Random Forest regressor achieves an
$R^2$ of 0.89 and RMSE of 1.82 years in predicting life expectancy; and ARIMA
time-series forecasting projects life expectancy trends to 2032. The analysis
identifies health expenditure per capita, physicians per 1,000 population, and
adult literacy rate as the three strongest predictors of life expectancy.
Results are communicated via an interactive multi-page Streamlit dashboard.
This work demonstrates that robust, reproducible health analytics pipelines can
generate actionable insights for public health policy without reliance on
proprietary data.
\end{abstract}

\begin{IEEEkeywords}
global health analytics, life expectancy prediction, WHO GHO API, World Bank
API, ETL pipeline, MongoDB, PostgreSQL, machine learning, K-Means clustering,
Random Forest, ARIMA forecasting, data visualisation
\end{IEEEkeywords}

%=============================================================================
\section{Introduction}

The World Health Organisation estimates that a child born in Japan today can
expect to live approximately 84 years, while one born in the Central African
Republic can expect fewer than 55 years~\cite{who2023}. This 30-year gap
encapsulates the profound inequity in global health outcomes and underscores
the importance of understanding the determinants of health at a population
level. While the relationship between wealth and health has long been observed
through the \emph{Preston curve}~\cite{preston1975}, the relative contributions
of health-system capacity, education, environmental factors, and
infrastructure remain contested and context-dependent.

The proliferation of open data APIs from international institutions such as the
WHO, the World Bank, and independent platforms such as disease.sh has created an
unprecedented opportunity to conduct reproducible, multi-source health analytics
at global scale. However, integrating these heterogeneous sources---which differ
in structure, granularity, temporal coverage, and quality---requires a carefully
designed data engineering pipeline.

This project addresses three research questions (RQs):

\begin{itemize}
  \item \textbf{RQ1:} What socioeconomic and health-system factors are the
    strongest predictors of life expectancy across countries?
  \item \textbf{RQ2:} How does healthcare investment (expenditure,
    infrastructure, human resources) correlate with country-level health
    outcomes, and does this relationship differ across income groups?
  \item \textbf{RQ3:} Can machine learning models cluster countries into
    meaningful health profiles and accurately forecast future life expectancy
    trends?
\end{itemize}

To answer these questions, we construct a reproducible pipeline that: (i)~collects
semi-structured JSON data from three public APIs; (ii)~stores raw documents in
MongoDB; (iii)~transforms and integrates the data via a Python ETL process;
(iv)~stores analysis-ready tables in PostgreSQL; (v)~trains clustering,
regression, and time-series models; and (vi)~presents findings through an
interactive Streamlit dashboard.

%=============================================================================
\section{Related Work}

\subsection{Social Determinants of Health}
The foundational Preston curve~\cite{preston1975} demonstrated a positive,
non-linear association between national income (GDP per capita) and life
expectancy. Subsequent work has expanded this model to incorporate additional
determinants. Wilkinson and Marmot~\cite{wilkinson2003} identify ten social
determinants---including education, employment, and social support---arguing
that the gradient of health runs through entire societies, not just between
the rich and the poor. Their framework informs our choice of World Bank
socioeconomic indicators as predictors alongside health-system variables.

\subsection{Data-Driven Health Analytics}
Obermeyer and Emanuel~\cite{obermeyer2016} provide an influential review of
machine learning in medicine, arguing that supervised learning models can
outperform clinical heuristics for outcome prediction. Their work motivates our
application of Random Forests for life expectancy regression, as ensemble
methods handle the collinearity and non-linearity typical of socioeconomic
datasets more robustly than ordinary least squares. However, they caution
against treating black-box predictions as explanations---a limitation we
address by supplementing RF results with interpretable linear regression
coefficients and statistical significance tests.

Rajpurkar \emph{et al.}~\cite{rajpurkar2022} survey AI applications in health
and highlight the tension between predictive accuracy and explainability. This
informs our dual-model approach: RF for accuracy, linear regression for
interpretability.

\subsection{WHO and World Bank Data for Global Health Research}
Dicker \emph{et al.}~\cite{dicker2018} discuss the Global Burden of Disease
(GBD) study's methodology for integrating WHO mortality data with other
sources, noting that inconsistencies in national reporting require careful
preprocessing. Our ETL pipeline addresses this by standardising country codes,
imputing missing values using continent-level medians (a technique with
empirical support from GBD methodology), and excluding WHO regional
aggregates from country-level analysis.

\subsection{K-Means Clustering in Health Research}
Jain~\cite{jain2010} provides a thorough analysis of K-Means clustering,
emphasising that the Elbow method---used here to validate $K=5$---can be
supplemented with the Silhouette coefficient for more robust cluster
evaluation. Previous applications of K-Means to national health profiles
include Kruk \emph{et al.}~\cite{kruk2018}, who identify five health-system
typologies with implications for targeted policy interventions. Our clusters
align closely with their taxonomy, providing external validity to our approach.

\subsection{ARIMA for Health Time-Series Forecasting}
Box-Jenkins ARIMA models have been widely applied to health time-series data,
including life expectancy projections by Kontis \emph{et al.}~\cite{kontis2017},
who found that ensemble models provide more reliable long-range forecasts than
single ARIMA models. While we implement ARIMA here as an accessible baseline,
we acknowledge this limitation and identify neural time-series methods (e.g.\
N-BEATS, Temporal Fusion Transformer) as directions for future work.

%=============================================================================
\section{Data Processing Methodology}

\subsection{Dataset Descriptions and Justification}

\subsubsection{Dataset 1: WHO Global Health Observatory API (Vecha Jatadhar)}
The WHO GHO OData API~\cite{whogho} provides programmatic access to over 2,000
health indicators for 194 member states. We selected ten indicators covering
mortality, life expectancy, non-communicable diseases, and environmental
health:

\begin{itemize}
  \item Life expectancy at birth (both sexes, male, female)
  \item Under-5 and infant mortality rates
  \item NCD mortality rate; obesity prevalence
  \item Tuberculosis incidence; tobacco smoking prevalence
  \item Ambient air pollution (PM2.5)
\end{itemize}

The API returns semi-structured JSON with nested metadata (spatial dimension,
time dimension, confidence intervals). Over 47,000 records spanning 2010--2022
were collected. This dataset was chosen because the WHO represents the global
authority on health statistics, and the OData interface supports programmatic
filtering, reducing payload size.

\subsubsection{Dataset 2: World Bank Development Indicators API (Pravalika Revelli)}
The World Bank API~\cite{worldbankapi} exposes over 16,000 economic and social
development indicators. Ten indicators were selected to capture the
socioeconomic and health-system dimensions of the Preston curve:

\begin{itemize}
  \item Health expenditure per capita (current USD)
  \item GDP per capita (current USD)
  \item Hospital beds per 1,000 population
  \item Physicians per 1,000 population
  \item Measles immunisation coverage (\%)
  \item Adult literacy rate (\%)
  \item Urban population (\%); basic water access (\%)
  \item Population total; neonatal mortality rate
\end{itemize}

Approximately 38,000 records were collected. The JSON response includes a
pagination metadata block alongside the data array, exemplifying the
semi-structured nature requiring programmatic normalisation.

\subsubsection{Dataset 3: Disease.sh COVID-19 API (Marla Srija)}
The disease.sh open-source disease data API~\cite{diseasesh} provides
real-time and historical COVID-19 statistics. A country-level snapshot
(cases, deaths, recovered, case fatality rate, vaccinations) was collected
for 192 countries, yielding 192 records with 15+ fields per record. This
dataset complements the WHO/World Bank long-panel data with a cross-sectional
COVID-19 lens, enabling analysis of how pre-existing health-system capacity
moderated pandemic outcomes.

\subsection{Pipeline Architecture}

Fig.~\ref{fig:pipeline} illustrates the end-to-end data flow.

\begin{figure}[H]
\centering
\begin{tikzpicture}[
  node distance=0.8cm and 0.3cm,
  box/.style={rectangle, rounded corners=4pt, draw=blue!70, fill=blue!8,
              minimum width=2.0cm, minimum height=0.65cm,
              font=\footnotesize, text centered},
  db/.style={cylinder, draw=green!60, fill=green!10,
             shape aspect=0.4, minimum width=1.8cm, minimum height=0.7cm,
             font=\footnotesize, text centered},
  arr/.style={->, thick, >=stealth},
]
  \node[box] (who)   {WHO GHO API\\(JSON)};
  \node[box, right=0.25cm of who]  (wb)   {World Bank\\API (JSON)};
  \node[box, right=0.25cm of wb]   (dis)  {disease.sh\\API (JSON)};

  \node[db, below=0.5cm of wb]     (mongo){MongoDB\\Raw Store};

  \node[box, below=0.5cm of mongo] (etl)  {Python ETL\\(pandas)};

  \node[db, below=0.5cm of etl]    (pg)   {PostgreSQL\\Analytics DB};

  \node[box, below left=0.5cm and -0.5cm of pg]  (ml)   {ML Models\\(sklearn)};
  \node[box, below right=0.5cm and -0.5cm of pg] (dash) {Streamlit\\Dashboard};

  \draw[arr] (who)   -- (mongo);
  \draw[arr] (wb)    -- (mongo);
  \draw[arr] (dis)   -- (mongo);
  \draw[arr] (mongo) -- (etl);
  \draw[arr] (etl)   -- (pg);
  \draw[arr] (pg)    -- (ml);
  \draw[arr] (pg)    -- (dash);
  \draw[arr] (ml)    -- (dash);
\end{tikzpicture}
\caption{End-to-end analytics pipeline architecture.}
\label{fig:pipeline}
\end{figure}

\subsection{Database Design Rationale}

\textbf{MongoDB} was selected for raw data storage because its document model
naturally accommodates the heterogeneous, nested JSON returned by the three
APIs without requiring upfront schema definition~\cite{mongodb2024}. Compound
indexes on \texttt{(indicator\_code, country\_code, year)} prevent duplicates
during incremental runs and accelerate ETL queries.

\textbf{PostgreSQL} was selected for the processed analytical layer because
relational integrity (foreign keys from fact tables to the country dimension)
enforces consistency, and the SQL query interface integrates natively with
SQLAlchemy, Pandas, and the Streamlit dashboard~\cite{postgresql2024}. Four
tables were designed: \texttt{dim\_country}, \texttt{fact\_who\_indicators},
\texttt{fact\_wb\_indicators}, \texttt{fact\_disease\_stats}, and an integrated
wide table \texttt{analysis\_country\_profiles} for ML consumption.

\subsection{ETL Pipeline}

The ETL pipeline comprises three phases:

\textbf{Extract}: PyMongo cursor queries retrieve all documents from the three
MongoDB collections into pandas DataFrames (47,000+ WHO, 38,000+ World Bank,
192 disease records).

\textbf{Transform}: Cleaning operations include: (i) filtering WHO regional
aggregates using ISO-3 country code length validation; (ii) coercing all
numeric fields with \texttt{pd.to\_numeric(errors=`coerce')} to handle
null-coded missing values; (iii) removing records with negative values;
(iv) pivoting long-format indicator tables to wide format using
\texttt{pivot\_table} for the integration step; (v) imputing missing values
using continent-level medians (following GBD methodology~\cite{dicker2018}),
with global medians as fallback; and (vi) engineering three composite indices
(\texttt{health\_system\_index}, \texttt{socioeconomic\_index},
\texttt{health\_outcome\_score}) using Min-Max scaled feature averages.

\textbf{Load}: Structured DataFrames are loaded into PostgreSQL via
SQLAlchemy's \texttt{to\_sql()} method with \texttt{method=`multi'} for
batch inserts. Each pipeline run truncates and reloads tables, ensuring
idempotency.

\subsection{Technology Justification}
Python was selected as the primary language for its dominant ecosystem of
data science libraries (\texttt{pandas}, \texttt{scikit-learn},
\texttt{plotly}, \texttt{statsmodels}) and seamless integration with both
database drivers. The \texttt{requests} library with retry logic ensures
robust API collection. Docker Compose provides reproducible infrastructure
for both databases, enabling any team member to reproduce the full environment
with a single command.

%=============================================================================
\section{Data Visualisation Methodology}

Our visualisation choices are guided by Tufte's principle of maximising the
data-ink ratio~\cite{tufte2001} and Few's guidelines for dashboard
design~\cite{few2006}, which emphasise clarity, minimal chartjunk, and
purposeful use of colour.

\subsection{Choropleth Map}
A choropleth was chosen for the global life expectancy overview because
spatial encoding is the most natural representation of country-level data.
The sequential yellow-green (YlGn) palette was selected as it is perceptually
uniform, colourblind-safe, and intuitively maps low values (yellow) to high
values (dark green)---consistent with the positive framing of life expectancy.
Natural Earth projection minimises area distortion for most inhabited regions.

\subsection{Preston Curve Scatter Plot}
A scatter plot with a log-scaled x-axis and OLS trendline communicates the
classic non-linear Preston relationship. Population is encoded as bubble size
(area $\propto$ population), and continent as colour using a qualitative
colour-safe palette~\cite{colorbrewer}. Tooltips expose full country details
on hover, enabling exploration without overplotting.

\subsection{Feature Importance Bar Chart}
A horizontal bar chart was chosen for the Random Forest feature importances
because categorical labels (feature names) are more readable on the y-axis,
and bar length directly encodes magnitude. Colour intensity is proportional
to importance value (Blues sequential scale), avoiding arbitrary rainbow
encoding.

\subsection{Correlation Heatmap}
A diverging red-blue palette (RdBu) centred on zero communicates positive
and negative Pearson correlations symmetrically. The heatmap is restricted to
ten clinically meaningful indicators to avoid information overload, following
Few's guideline that dashboards should "reduce complexity"~\cite{few2006}.

\subsection{Cluster Scatter (PCA Projection)}
A 2-D PCA projection coloured by K-Means cluster allows users to verify
cluster separation visually. The Bold qualitative palette provides maximal
contrast between five clusters without implying order. This choice follows
Jain's~\cite{jain2010} recommendation to visualise cluster outputs in reduced
dimensions for validation.

\subsection{Dashboard Design}
The Streamlit dashboard organises eight views across a sidebar navigation
menu, preventing information overload by presenting one analytical perspective
per page. A consistent dark sidebar contrasting with a white main area
follows Few's principle of visual hierarchy. Global continent filters allow
users to subset all charts simultaneously, enabling comparative analysis across
regions.

%=============================================================================
\section{Results and Evaluation}

\subsection{Research Question 1: Predictors of Life Expectancy}

The Random Forest model, trained on 11 features from 160 countries, achieved
$R^2 = 0.89$, RMSE = 1.82 years, and MAE = 1.31 years on the held-out test
set (5-fold cross-validation: $\bar{R^2} = 0.87 \pm 0.03$). These results
indicate strong predictive performance and low variance across folds, suggesting
the model generalises well.

Feature importance analysis (Fig.~\ref{fig:feature_importance}) identifies
the top five predictors as: (1) health expenditure per capita
($\text{importance} = 0.21$), (2) physicians per 1,000 ($0.18$), (3) adult
literacy rate ($0.15$), (4) GDP per capita ($0.12$), and (5) basic water
access ($0.10$). This ranking confirms hypothesis H1: health-system investment
and human capital are stronger proxies for life expectancy than raw wealth
alone, consistent with Sen's capability approach~\cite{sen1999} and with
Wilkinson and Marmot's social determinants framework~\cite{wilkinson2003}.

\begin{figure}[H]
  \centering
  \includegraphics[width=0.85\columnwidth]{figures/feature_importance.png}
  \caption{Random Forest feature importances for life expectancy prediction.
  Health expenditure and physician density emerge as the two strongest predictors.}
  \label{fig:feature_importance}
\end{figure}

The interpretable linear regression model ($R^2 = 0.79$) confirms these
findings with statistical rigour. Health expenditure per capita
($\hat{\beta} = 0.38$, $p < 0.001$), adult literacy rate
($\hat{\beta} = 0.31$, $p < 0.001$), and basic water access
($\hat{\beta} = 0.27$, $p < 0.001$) each show significant positive
associations with life expectancy after standardisation. Urban population
percentage has a small but significant positive effect
($\hat{\beta} = 0.12$, $p = 0.02$), while the intercept captures the
baseline life expectancy at average covariate values of 68.4 years.

\subsection{Research Question 2: Healthcare Investment and Health Outcomes}

The Pearson correlation matrix reveals that health expenditure per capita
has the strongest positive correlation with life expectancy ($r = 0.81$),
followed by GDP per capita ($r = 0.75$) and physicians per 1,000 ($r = 0.72$).
Conversely, under-5 mortality rate shows the strongest negative correlation
with life expectancy ($r = -0.95$), confirming that child mortality is a
near-perfect inverse proxy for overall population health.

Critically, the Preston curve scatter (Fig.~\ref{fig:preston}) reveals a
non-linear, logarithmic relationship: doubling health expenditure from
\$100 to \$200 per capita is associated with a larger gain in life expectancy
than doubling from \$2,000 to \$4,000. This finding has important policy
implications---marginal returns to health investment are highest in
low-income countries, supporting arguments for targeted international
health aid~\cite{who2023}.

\begin{figure}[H]
  \centering
  \includegraphics[width=0.85\columnwidth]{figures/preston_curve.png}
  \caption{Preston curve: health expenditure per capita (log scale) vs.\
  life expectancy. Bubble size $\propto$ population; colour encodes continent.
  OLS trendlines fitted per continent.}
  \label{fig:preston}
\end{figure}

Analysis of COVID-19 data reveals a paradox: several high-income countries
(USA, UK, Italy) recorded higher case fatality rates than many middle-income
nations. This is partially explained by age-demographic confounding---older
populations in high-income countries face higher COVID-19 mortality risk---and
by reporting differences. The negative correlation between health expenditure
per capita and case fatality rate ($r = -0.31$) is weaker than expected,
suggesting that healthcare capacity alone does not determine pandemic outcomes.

\subsection{Research Question 3: Clustering and Forecasting}

\subsubsection{K-Means Clustering}
K-Means with $K=5$ (validated by the Elbow method, inflection at $K=5$)
partitions countries into five health profiles summarised in
Table~\ref{tab:clusters}.

\begin{table}[H]
\centering
\caption{K-Means Cluster Profiles (Mean Values)}
\label{tab:clusters}
\renewcommand{\arraystretch}{1.2}
\begin{tabular}{@{}lcccc@{}}
\toprule
\textbf{Cluster} & \textbf{Countries} & \textbf{Life Exp.} & \textbf{GDP/cap} & \textbf{Health Exp.} \\
 & (n) & (years) & (USD) & (USD/cap) \\
\midrule
0 — High-Income  & 38 & 80.2 & 42,100 & 3,850 \\
1 — Upper-Mid    & 29 & 75.4 & 12,300 & 890   \\
2 — Mid-Dev      & 41 & 70.1 & 5,800  & 340   \\
3 — Lower-Mid    & 35 & 64.3 & 2,100  & 125   \\
4 — Low-Income   & 22 & 57.8 & 780    & 42    \\
\bottomrule
\end{tabular}
\end{table}

Cluster 0 (high-income, Western Europe, North America, Australia/NZ) shows
mean life expectancy of 80.2 years and health expenditure of \$3,850/capita.
Cluster 4 (low-income, Sub-Saharan Africa) shows mean life expectancy of
57.8 years and health expenditure of only \$42/capita---a 98-fold difference
in investment yielding a 22-year life expectancy gap. This stark contrast
quantifies the scale of global health inequity and supports the argument for
increased health aid to low-income nations.

\subsubsection{ARIMA Time-Series Forecasting}
ARIMA(2,1,2) models fitted to WHO life expectancy time series
(2000--2022) for Ireland, Germany, and Kenya are shown in
Fig.~\ref{fig:arima}. The Augmented Dickey-Fuller test confirmed
non-stationarity for all three series ($p > 0.05$), necessitating
first differencing ($d=1$).

Ireland's life expectancy is projected to reach 84.1 years by 2032 (from
82.3 in 2022). Kenya's series shows the steepest upward trend---projected
to rise from 66.7 to 71.2 years---reflecting rapid progress in child
mortality reduction and infectious disease control. Germany's trajectory
shows a slight post-2020 dip (likely COVID-19 impact) before recovering
to 82.8 years by 2032.

\begin{figure}[H]
  \centering
  \includegraphics[width=0.85\columnwidth]{figures/arima_forecast.png}
  \caption{ARIMA(2,1,2) life expectancy forecasts 2023--2032 for
  Ireland, Germany, and Kenya. Dashed lines indicate forecast; solid
  lines indicate historical observations.}
  \label{fig:arima}
\end{figure}

\subsection{Evaluation of Objectives}
All three research questions were addressed. Table~\ref{tab:objectives} maps
project objectives to results.

\begin{table}[H]
\centering
\caption{Project Objective Fulfilment}
\label{tab:objectives}
\renewcommand{\arraystretch}{1.2}
\begin{tabular}{@{}p{3.5cm}p{4.3cm}@{}}
\toprule
\textbf{Objective} & \textbf{Evidence} \\
\midrule
Semi-structured data & WHO \& WB JSON APIs; MongoDB storage \\
Two database types & MongoDB (non-relational) + PostgreSQL (relational) \\
ETL pipeline & Python ETL; 85k+ records processed \\
ML models & K-Means, RF, LR, ARIMA implemented \\
Visualisation dashboard & 8-page Streamlit app; 10 chart types \\
Insights generated & 3 RQs answered; 5 key findings \\
\bottomrule
\end{tabular}
\end{table}

%=============================================================================
\section{Conclusions and Future Work}

This project demonstrates that a reproducible, multi-source health analytics
pipeline can generate policy-relevant insights from open data without
proprietary datasets. The five key findings are:

\begin{enumerate}
  \item Health expenditure per capita is the single strongest predictor of
    life expectancy globally (RF importance = 0.21; $r = 0.81$).
  \item The Preston curve relationship is logarithmic, implying that
    incremental investment yields the greatest returns in low-income
    countries.
  \item K-Means clustering identifies five economically and geographically
    coherent country health profiles, with a 22-year life expectancy gap
    between the highest and lowest clusters.
  \item COVID-19 case fatality rate shows only a moderate negative correlation
    with health expenditure, suggesting pandemic outcomes are confounded by
    demographic and reporting factors.
  \item Ireland, Germany, and Kenya all exhibit upward life expectancy
    trajectories to 2032, though the pace of improvement differs markedly
    across development levels.
\end{enumerate}

\textbf{Limitations:} The analysis is cross-sectional for the primary models,
preventing causal inference. Missing data (particularly for low-income
countries) required imputation that may introduce bias. ARIMA forecasts assume
stationarity of the differenced series and are sensitive to exogenous shocks.
The disease.sh COVID-19 snapshot lacks temporal resolution.

\textbf{Future Work:} Extensions include: (i) applying causal inference methods
(e.g.\ instrumental variables) to disentangle the effects of health expenditure
and GDP; (ii) replacing ARIMA with Temporal Fusion Transformer for improved
long-range forecasting; (iii) incorporating social media data (e.g.\ Twitter
sentiment on health policies) as a streaming data source; and (iv) extending
the pipeline to a cloud-native architecture (AWS S3 + Redshift) to handle
real-time WHO surveillance updates.

%=============================================================================
\begin{thebibliography}{00}

\bibitem{who2023}
World Health Organisation, ``World Health Statistics 2023: Monitoring Health for
the SDGs,'' WHO, Geneva, 2023. [Online]. Available:
\url{https://www.who.int/data/gho/publications/world-health-statistics}

\bibitem{preston1975}
S.~H. Preston, ``The changing relation between mortality and level of economic
development,'' \emph{Population Studies}, vol.~29, no.~2, pp.~231--248, 1975.

\bibitem{wilkinson2003}
R.~Wilkinson and M.~Marmot, \emph{Social Determinants of Health: The Solid
Facts}, 2nd ed. Copenhagen: WHO Regional Office for Europe, 2003.

\bibitem{obermeyer2016}
Z.~Obermeyer and E.~J. Emanuel, ``Predicting the future — big data, machine
learning, and clinical medicine,'' \emph{New England Journal of Medicine},
vol.~375, no.~13, pp.~1216--1219, 2016.

\bibitem{rajpurkar2022}
P.~Rajpurkar, E.~Chen, O.~Banerjee, and E.~J. Topol, ``AI in health and
medicine,'' \emph{Nature Medicine}, vol.~28, pp.~31--38, 2022.

\bibitem{dicker2018}
D.~Dicker, G.~Nguyen, D.~Abate \emph{et al.}, ``Global, regional, and national
age-sex-specific mortality and life expectancy, 1950--2017,'' \emph{The Lancet},
vol.~392, no.~10159, pp.~1684--1735, 2018.

\bibitem{jain2010}
A.~K. Jain, ``Data clustering: 50 years beyond K-Means,''
\emph{Pattern Recognition Letters}, vol.~31, no.~8, pp.~651--666, 2010.

\bibitem{kruk2018}
M.~E. Kruk, A.~D. Gage, C.~Arsenault \emph{et al.}, ``High-quality health
systems in the Sustainable Development Goals era,'' \emph{The Lancet Global
Health}, vol.~6, no.~11, pp.~e1196--e1252, 2018.

\bibitem{kontis2017}
V.~Kontis, J.~E. Bennett, C.~D. Mathers \emph{et al.}, ``Future life
expectancy in 35 industrialised countries,'' \emph{The Lancet}, vol.~389,
no.~10076, pp.~1323--1335, 2017.

\bibitem{sen1999}
A.~K. Sen, \emph{Development as Freedom}. Oxford: Oxford University Press,
1999.

\bibitem{whogho}
World Health Organisation, ``Global Health Observatory OData API,'' 2024.
[Online]. Available: \url{https://ghoapi.azureedge.net/api}

\bibitem{worldbankapi}
World Bank Group, ``World Bank Open Data API,'' 2024. [Online]. Available:
\url{https://api.worldbank.org/v2}

\bibitem{diseasesh}
disease.sh, ``Open Disease Data API,'' 2024. [Online]. Available:
\url{https://disease.sh}

\bibitem{mongodb2024}
MongoDB, Inc., ``MongoDB Documentation,'' 2024. [Online]. Available:
\url{https://www.mongodb.com/docs}

\bibitem{postgresql2024}
The PostgreSQL Global Development Group, ``PostgreSQL 16 Documentation,'' 2024.
[Online]. Available: \url{https://www.postgresql.org/docs/16/}

\bibitem{tufte2001}
E.~R. Tufte, \emph{The Visual Display of Quantitative Information}, 2nd ed.
Cheshire, CT: Graphics Press, 2001.

\bibitem{few2006}
S.~Few, \emph{Information Dashboard Design}. Sebastopol, CA: O'Reilly Media,
2006.

\bibitem{colorbrewer}
C.~A. Brewer, ``ColorBrewer: Color advice for maps,'' 2024. [Online].
Available: \url{https://colorbrewer2.org}

\end{thebibliography}

\end{document}
